\documentclass[11pt,a4paper]{article}
\usepackage[T1]{fontenc}
\usepackage{lmodern,amsmath,amssymb,amsthm,microtype,booktabs}
\usepackage[margin=25mm]{geometry}
\usepackage[colorlinks=true,urlcolor=blue,linkcolor=blue]{hyperref}
\newcommand{\dd}{\mathop{}\!\mathrm d}
\newcommand{\sech}{\operatorname{sech}}
\newcommand{\csch}{\operatorname{csch}}
\newcommand{\HH}{\mathcal H}
\title{The observer mass spectrum in dS$_3$:\\spacetime normalization and the BD projector}
\author{Alan Day \and Chengheng\thanks{AI research assistant (OpenAI).}}
\date{Working research note --- 31 October 2026}
\begin{document}\maketitle
\begin{abstract}
We distinguish fixed-mass Klein--Gordon normalization from the mass-delta
normalization induced by spacetime $L^2$ in a specified $G=0$ observer model.
A Darboux polar decomposition gives the full radial spectrum, including
its two continuum channels and bound states. The Bunch--Davies line and
its antipodal image have a nonzero spacetime overlap in dS$_3$.
We obtain the orthogonal observer projector and the resulting two-leg
rescaling of the previously computed local vertex. Bounded compressed
insertions provide a controlled intermediate-state identity.
Group averaging, the physical state, and finite-gravity sewing remain open.
\end{abstract}

\section{Spacetime $L^2$ and a complete radial resolution}
Set the radius to one, use $G=0$, and normalize $\int |Y_{lm}|^2\dd\Omega=1$.
For
\[
 ds^2=-\dd t^2+\cosh^2t\,\dd\Omega_2^2,\qquad
 F=\sum_{lm}f_{lm}(t)Y_{lm},\qquad \chi_{lm}=\cosh t\,f_{lm},
\]
the map $R:F\mapsto(\chi_{lm})$ is unitary from spacetime $L^2$
to $\bigoplus_{lm}L^2(\mathbb R,\dd t)$ and
\begin{equation}
 R\Box R^{-1}=\bigoplus_{lm}(1+H_l),\qquad
 H_l=-\partial_t^2-l(l+1)\sech^2t .
\end{equation}
Each $H_l$ is self-adjoint on $H^2(\mathbb R)$ with core
$C_c^\infty(\mathbb R)$, since its potential is real and bounded.
The full direct-sum domain additionally requires
$\sum_{lm}\|(1+H_l)\chi_{lm}\|^2<\infty$.
The finite angular-mode core fixes the closure; the larger invariant
symmetric domain $C_c^\infty(\mathrm{dS}_3)$ has the same closure.
It follows that the Casimir spectral projections commute with proper,
orthochronous isometries.

The reflectionless potential and its bound states are standard
\cite{LMW}. To fix completeness and normalization explicitly, let
$D_j=\partial_t-j\tanh t$. Then
\begin{equation}
 D_j^\dagger D_j=H_{j-1}+j^2,\quad
 D_jD_j^\dagger=H_j+j^2,\quad
 W_j=D_j(H_{j-1}+j^2)^{-1/2}.
\end{equation}
Inductively $H_{j-1}\geq-(j-1)^2$, so the inverse square root is bounded.
$W_j$ is an isometry onto $(\ker D_j^\dagger)^\perp$, intertwines the two
Hamiltonians, and the omitted kernel is spanned by $\sech^jt$ at energy
$-j^2$. Starting with the ordinary Fourier transform of $H_0$ proves:
\begin{itemize}
\item the positive spectrum $\nu^2>0$ is purely absolutely continuous
with two channels;
\item each radial $H_l$ has the simple eigenvalues $-k^2$, $1\leq k\leq l$;
\item the bounded zero-energy solution $P_l(\tanh t)$ is a resonance,
not an $L^2$ eigenvector or a spectral atom.
\end{itemize}
Thus the full $\Box$ has point eigenvalues $1-k^2$, each with infinitely
many angular modes $l\geq k$. We do not identify these with a group
discrete series without a separate representation analysis.

Let $d_l(\nu)=\prod_{j=1}^l\sqrt{\nu^2+j^2}$. A delta-normalized
continuum basis, with $\nu>0$ for each sign, is
\begin{equation}
 e_{\nu l}^{\mp}(t)=
 \frac{i^lD_l\cdots D_1e^{\mp i\nu t}}{\sqrt{2\pi}\,d_l(\nu)}.
\end{equation}
A real normalized bound basis is
\begin{equation}
 b_{lk}(t)=
 \frac{D_l\cdots D_{k+1}\sech^kt}
 {\sqrt{B(1/2,k)}\prod_{j=k+1}^l\sqrt{j^2-k^2}}.
\end{equation}
For $a_\pm(\nu)=\langle e^\pm_{\nu l},\chi\rangle$ and
$c_k=\langle b_{lk},\chi\rangle$, the exact resolution is
\begin{align}
 \|\chi\|^2&=\int_0^\infty(|a_-|^2+|a_+|^2)\dd\nu+
                 \sum_{k=1}^l|c_k|^2,\\
 \chi&=\int_0^\infty(e^-a_-+e^+a_+)\dd\nu+\sum_{k=1}^lb_{lk}c_k .
\end{align}
The inverse is an $L^2$ limit. A single continuum channel is not a
resolution of the full position delta function.

\section{The BD and antipodal Gram matrix}
Use the KG-normalized BD mode $u_{\nu lm}=f_{\nu l}Y_{lm}$ from the
preceding note \cite{local}, consistent with \cite{MM}. Set
\[
 q=e^{-\pi\nu},\qquad N=[2\nu(1-q^2)]^{-1/2},\qquad
 \chi_B=\cosh t\,f_{\nu l}=\sqrt{2\pi}N(e^- -qe^+).
\]
The antipodal pullback $A:(t,n)\mapsto(-t,-n)$ is a self-adjoint
unitary on spacetime $L^2$, commuting with the proper isometry group.
Keeping the angular harmonic $Y_{lm}$ fixed, its radial mode is
\begin{equation}
 \chi_A(t)=(-1)^l\chi_B(-t)=(-1)^l\chi_B(t)^*
             =\sqrt{2\pi}N(e^+ -qe^-).
\end{equation}
Complex conjugation of the full spherical harmonic also reverses $m$;
the radial and full conjugations must be distinguished.

The continuum resolution now gives
\begin{equation}
 \langle u_{\nu lm}^{r},u_{\nu'l'm'}^{s}\rangle_{L^2}
 =G_{rs}(\nu)\delta(\nu-\nu')\delta_{ll'}\delta_{mm'},\qquad
 G=\frac{\pi}{\nu}
 \begin{pmatrix}\coth\pi\nu&-\csch\pi\nu\\
                 -\csch\pi\nu&\coth\pi\nu\end{pmatrix}.
 \label{eq:gram}
\end{equation}
Its eigenvalues are $(\pi/\nu)\tanh(\pi\nu/2)$ and
$(\pi/\nu)\coth(\pi\nu/2)$, with $\det G=(\pi/\nu)^2$.
This positive spacetime Gram matrix is different from the fixed-mass
KG signature $(+1,-1)$. In a basis using $\chi_B^*$ as the second radial
mode, the off-diagonal entry acquires $(-1)^l$.

KL's dS$_2$ construction \cite[eqs.\ (4.9)--(4.11)]{KL} has an
orthogonal antipodal pair. Equation \eqref{eq:gram} is a fresh
dS$_3$ derivation; that cross-channel orthogonality must not be
transplanted between dimensions. The agreement of the diagonal
normalizing function below does not imply agreement of the two Gram
matrices, and is not a claim of an error in the dS$_2$ result.

If both nonorthogonal channels are used, the full continuous projector
requires their inverse Gram matrix:
\[
 P_{\mathrm{ac}}=\sum_{lm}\int_0^\infty
 |u^r_{\nu lm}\rangle (G^{-1})_{rs}
 \langle u^s_{\nu lm}|\dd\nu .
\]
Separately normalizing each vector and adding its rank-one projector
does not give the identity.

\section{An explicit observer subspace and its measure}
Choose and retain a mass cutoff $\nu_*>0$. Define
\begin{equation}
 \mu(\nu)=\frac{\nu\tanh\pi\nu}{\pi},\qquad
 v_{\nu lm}=\sqrt{\mu(\nu)}\,u_{\nu lm},\qquad
 (Ja)(X)=\sum_{lm}\int_{\nu_*}^\infty v_{\nu lm}(X)a_{lm}(\nu)\dd\nu.
 \label{eq:J}
\end{equation}
Then $J$ is an isometric embedding of
$\HH_{\mathrm{obs}}=\int_{\nu_*}^\infty{}^\oplus\ell^2(l,m)\dd\nu$
into spacetime $L^2$:
$\|Ja\|^2=\int_{\nu_*}^\infty\sum_{lm}|a_{lm}|^2\dd\nu$.
In the ordered $(e^-,e^+)$ continuum basis,
\begin{equation}
 v_B=\frac{(1,-q)^T}{\sqrt{1+q^2}},\quad
 v_\perp=\frac{(q,1)^T}{\sqrt{1+q^2}},\quad
 P_B=\frac1{1+q^2}\begin{pmatrix}1&-q\\-q&q^2\end{pmatrix}.
\end{equation}
The observer projector is $1_{\nu\geq\nu_*}P_B$ on the continuum and
zero on every bound mode. Consequently $J^\dagger J=1$ and
$JJ^\dagger=P_{\mathrm{obs}}$, a norm-one orthogonal projection.

Both continuum channels have the same adjacent-angular-momentum
boost coefficient
\begin{equation}
 \beta_{j,m}(\nu)=
 \sqrt{\frac{(j^2-m^2)(j^2+\nu^2)}{(2j-1)(2j+1)}} .
\end{equation}
This follows from the Darboux and spherical-harmonic recurrences, and
agrees with \cite[eqs.\ (20)--(21)]{MM} at $d=2$, $\sigma=-1+i\nu$.
Rotations also act identically. The matrix $P_B$ depends on $\nu$ only;
hence the projection commutes with proper isometries. Equivalently, the
BD space is invariant, its antipodal image is invariant, and these two
lines span the continuum multiplicity space for $\nu>0$.

In contrast to an orthogonal complement, the antipodal unit vector obeys
\begin{equation}
 \langle v_B,Av_B\rangle=-\sech\pi\nu,\qquad
 v_\perp=\frac{Av_B+\sech(\pi\nu)v_B}{\tanh\pi\nu}.
 \label{eq:antipode}
\end{equation}
The condition number of the double Gram is $\coth^2(\pi\nu/2)$.
It diverges at zero, while $P_B$ remains bounded and $\mu\sim\nu^2$.
At large $\nu$, $\mu\sim\nu/\pi$ and the antipodal overlap is
exponentially small. These limits do not remove the declared cutoff.

Restoring the radius, $u_\ell=\ell^{-1/2}u_1$,
$G_\ell=\ell^2G$, and $v_\ell=\ell^{-3/2}v_1$.
For physical effective mass $m=\sqrt{1+\nu^2}/\ell$,
\begin{equation}
 \frac{\dd\nu}{\dd m}=\frac{\ell^2m}{\nu},\qquad
 v_m=\sqrt{\frac{\ell^2m}{\nu}}\,v_\nu .
\end{equation}
For $\lambda=1+\nu^2$, the delta-normalized mode is
$v_\lambda=v_\nu/\sqrt{2\nu}$ instead.
These are changes of spectral variable, not extra group Plancherel
densities to multiply into $\dd\nu$.

\section{What this fixes in the local insertion}
For observer slots $(o,i)$ and a KG-normalized free-field probe $p$, the
local coefficient of the preceding note becomes
\begin{equation}
 \boxed{\ C_\eta^{\mathrm{obs}}(\nu_o,\nu_i,\nu_p)
 =\sqrt{\mu(\nu_o)\mu(\nu_i)}\,C_\eta^{\mathrm{KG}}(\nu_o,\nu_i,\nu_p)\ } .
\end{equation}
Only the observer legs change normalization; the field's creation and
annihilation operators retain their KG convention.
The coefficient has dimension $\ell^{-1/2}$ with the chosen unit local
coupling. This fixes the legs of this kinematic realization. It does
not yet fix the earlier unknown physical two-point factor $n$, which
also involved a state, group reduction and sewing norms.

A bounded version makes the intermediate projector unambiguous.
For $g\in L^\infty(\mathrm{dS}_3)$, define
\begin{equation}
 A_g=J^\dagger M_gJ,\qquad \|A_g\|\leq\|g\|_\infty,\qquad
 A_g^\dagger=A_{\bar g}.
\end{equation}
This assignment is isometry-covariant, and
\begin{equation}
 A_gA_h=J^\dagger M_gP_{\mathrm{obs}}M_hJ,\qquad
 A_{gh}-A_gA_h=J^\dagger M_g(1-P_{\mathrm{obs}})M_hJ.
 \label{eq:leakage}
\end{equation}
For real $g=h$ the last expression is positive, being $B^\dagger B$
with $B=(1-P_{\mathrm{obs}})M_gJ$.
Replacing the intermediate spectral resolution by the full position
delta would therefore change even this bounded two-point quantity.

Each fixed free-field angular mode $u_p$ is bounded. Keeping a finite
set $S$ of such probe modes, let
\[
 O_S=\sum_{p\in S}\big(A_{u_p}\otimes a_p+
                         A_{u_p}^\dagger\otimes a_p^\dagger\big).
\]
Finite-particle vectors are analytic for this finite sum with bounded
coefficients, so it is essentially self-adjoint on that core.
Its vacuum two-point identity is
\begin{equation}
 \langle\psi,0|O_S^2|\psi,0\rangle
 =\sum_{p\in S}\|A_{u_p}^\dagger\psi\|^2
 =\sum_{p\in S,lm}\int_{\nu_*}^\infty
 |\langle\nu lm|A_{u_p}^\dagger|\psi\rangle|^2\dd\nu .
\end{equation}
A finite probe-mode truncation breaks full de Sitter invariance.
Removing it requires independent ultraviolet and domain estimates.
This kinematic identity is not a group-averaged Hartle--Hawking
two-point function or a finite-$G$ result.

\section{Independent checks and the next physical step}
The saved diagnostics contain 150 passed checks, with the two failed
initial finite-box targets retained separately. Matrix, mode and bound
state calculations use 50 decimal digits. Four $\sech^p t$ test functions
check Parseval by actually integrating the continuous spectral weight
and adding bound-state overlaps; independent time integrals check their
Fourier transforms. The angular sectors include $l=0,1,2,3$, so deleting
bound states produces a substantial, visible norm deficit.

A compact packet with support $0.6\leq\nu\leq1.8$ and normalized amplitude
$h(\nu)\propto(\nu-0.6)^4(1.8-\nu)^4$ independently checks
\eqref{eq:J} and \eqref{eq:antipode}. Mass quadrature uses 140/220 nodes
and time quadrature 360/540 nodes. At $T=20$, the finite-box norm deficit
is $3.67314\times10^{-6}$, exceeding the initial $2\times10^{-7}$
target; this failure was preserved. At $T=40,80$ the deficits are
$5.46933\times10^{-9}$ and $1.37811\times10^{-11}$.
The predicted antipodal overlap is approximately $-0.05055152806$;
it is not zero.

The missing tail estimate follows from Fourier Plancherel.
For $b_\pm=h(1,-q)_\pm/\sqrt{1+q^2}$, the first three endpoint derivatives
vanish, so
\begin{equation}
 \int_{|t|>T}|\chi(t)|^2\dd t
 \leq T^{-8}\int_{0.6}^{1.8}
       (|b_-^{(4)}|^2+|b_+^{(4)}|^2)\dd\nu .
\end{equation}
Cauchy--Schwarz bounds the omitted cross overlaps by the same coefficient.
Its 64/96-node evaluations give approximately $3.26727\times10^6$,
hence a conservative $T=20$ estimate $1.27628\times10^{-4}$.
The inequality is analytic; the evaluated coefficient and all finite
quadratures are empirical, not interval certificates.
The failed short-box tests were not erased or relabeled as a
high-precision delta-normalization test.

The next task is the two-observer group-averaged inner product and its
state, with an explicit Haar convention and reflection terms.
Only then can this fixed local normalization be tested against the
physical two-point spectrum. Independent free-probe disconnected
four-point limits, removal of matter truncations, finite-$G$ dynamics
and the full holographic objective remain required.

\begin{thebibliography}{9}
\bibitem{LMW} P.~Lagogiannis, A.~Maloney and Y.~Wang,
\emph{Odd-dimensional de Sitter Space is Transparent},
\href{https://arxiv.org/html/1106.2846v1}{arXiv:1106.2846v1}.
\bibitem{MM} D.~Marolf and I.~A.~Morrison,
\emph{Group Averaging for de Sitter free fields},
\href{https://arxiv.org/html/0810.5163v1}{arXiv:0810.5163v1}.
\bibitem{KL} D.~Kolchmeyer and H.~Liu,
\emph{Chaos and the Emergence of the Cosmological Horizon},
\href{https://arxiv.org/html/2411.08090v1}{arXiv:2411.08090v1}.
\bibitem{local} Alan Day and Chengheng,
\emph{Local scalar insertions in dS$_3$: from normalized BD modes
to global CG coefficients}, internal research note, 31 October 2026.
\end{thebibliography}
\end{document}
